% concatenated from Fin-dim_springer_24.06.2026.tex
%Version 3.1 December 2024
% See section 11 of the User Manual for version history
%
%%%%%%%%%%%%%%%%%%%%%%%%%%%%%%%%%%%%%%%%%%%%%%%%%%%%%%%%%%%%%%%%%%%%%%
%%                                                                 %%
%% Please do not use \input{...} to include other tex files.       %%
%% Submit your LaTeX manuscript as one .tex document.              %%
%%                                                                 %%
%% All additional figures and files should be attached             %%
%% separately and not embedded in the \TeX\ document itself.       %%
%%                                                                 %%
%%%%%%%%%%%%%%%%%%%%%%%%%%%%%%%%%%%%%%%%%%%%%%%%%%%%%%%%%%%%%%%%%%%%%

%%\documentclass[referee,sn-basic]{sn-jnl}% referee option is meant for double line spacing

%%=======================================================%%
%% to print line numbers in the margin use lineno option %%
%%=======================================================%%

%%\documentclass[lineno,pdflatex,sn-basic]{sn-jnl}% Basic Springer Nature Reference Style/Chemistry Reference Style

%%=========================================================================================%%
%% the documentclass is set to pdflatex as default. You can delete it if not appropriate.  %%
%%=========================================================================================%%

%%\documentclass[sn-basic]{sn-jnl}% Basic Springer Nature Reference Style/Chemistry Reference Style

%%Note: the following reference styles support Namedate and Numbered referencing. By default the style follows the most common style. To switch between the options you can add or remove “Numbered” in the optional parenthesis. 
%%The option is available for: sn-basic.bst, sn-chicago.bst%  

%%\documentclass[pdflatex,sn-nature]{sn-jnl}% Style for submissions to Nature Portfolio journals
%%\documentclass[pdflatex,sn-basic]{sn-jnl}% Basic Springer Nature Reference Style/Chemistry Reference Style
\documentclass[pdflatex,sn-mathphys-num]{sn-jnl}% Math and Physical Sciences Numbered Reference Style
%%\documentclass[pdflatex,sn-mathphys-ay]{sn-jnl}% Math and Physical Sciences Author Year Reference Style
%%\documentclass[pdflatex,sn-aps]{sn-jnl}% American Physical Society (APS) Reference Style
%%\documentclass[pdflatex,sn-vancouver-num]{sn-jnl}% Vancouver Numbered Reference Style
%%\documentclass[pdflatex,sn-vancouver-ay]{sn-jnl}% Vancouver Author Year Reference Style
%%\documentclass[pdflatex,sn-apa]{sn-jnl}% APA Reference Style
%%\documentclass[pdflatex,sn-chicago]{sn-jnl}% Chicago-based Humanities Reference Style

%%%% Standard Packages
%%<additional latex packages if required can be included here>

\usepackage{graphicx}%
\usepackage{multirow}%
\usepackage{amsmath,amssymb,amsfonts}%
\usepackage{amsthm}%
\usepackage{mathrsfs}%
\usepackage[title]{appendix}%
\usepackage{xcolor}%
\usepackage{textcomp}%
\usepackage{manyfoot}%
\usepackage{booktabs}%
\usepackage{algorithm}%
\usepackage{algorithmicx}%
\usepackage{algpseudocode}%
\usepackage{listings}%

\usepackage[T1]{fontenc}
\usepackage[polish,english]{babel}
\usepackage[utf8]{inputenc}

%%%%

%%%%%=============================================================================%%%%
%%%%  Remarks: This template is provided to aid authors with the preparation
%%%%  of original research articles intended for submission to journals published 
%%%%  by Springer Nature. The guidance has been prepared in partnership with 
%%%%  production teams to conform to Springer Nature technical requirements. 
%%%%  Editorial and presentation requirements differ among journal portfolios and 
%%%%  research disciplines. You may find sections in this template are irrelevant 
%%%%  to your work and are empowered to omit any such section if allowed by the 
%%%%  journal you intend to submit to. The submission guidelines and policies 
%%%%  of the journal take precedence. A detailed User Manual is available in the 
%%%%  template package for technical guidance.
%%%%%=============================================================================%%%%

%% as per the requirement new theorem styles can be included as shown below
\theoremstyle{thmstyleone}%
\newtheorem{theorem}{Theorem}%  meant for continuous numbers
%%\newtheorem{theorem}{Theorem}[section]% meant for sectionwise numbers
%% optional argument [theorem] produces theorem numbering sequence instead of independent numbers for Proposition
\newtheorem{proposition}[theorem]{Proposition}% 
%%\newtheorem{proposition}{Proposition}% to get separate numbers for theorem and proposition etc.

\newtheorem{lemma}[theorem]{Lemma}% 

\theoremstyle{thmstyletwo}%
\newtheorem{example}{Example}%
\newtheorem{remark}{Remark}%

\theoremstyle{thmstylethree}%
\newtheorem{definition}{Definition}%

%\newproof{pf}{proof}
%\newproof{pot}{Proof of Theorem \ref{thm2}}

\raggedbottom
%%\unnumbered% uncomment this for unnumbered level heads

\begin{document}
	
	\title[On closed linear subspaces embedded into Banach spaces
	and their finite-dimensionality]{On closed linear subspaces embedded into Banach spaces
		and their finite-dimensionality}
	
	%%=============================================================%%
	%% GivenName	-> \fnm{Joergen W.}
	%% Particle	-> \spfx{van der} -> surname prefix
	%% FamilyName	-> \sur{Ploeg}
	%% Suffix	-> \sfx{IV}
	%% \author*[1,2]{\fnm{Joergen W.} \spfx{van der} \sur{Ploeg} 
		%%  \sfx{IV}}\email{iauthor@gmail.com}
	%%=============================================================%%
	

	
	\author[1]{\fnm{Alexander} \sur{Balinsky}}\email{BalinskyA@cardiff.ac.uk}

	
		\author*[2,3]{\fnm{Yarema} \sur{Prykarpatskyy}}\email{yarema.prykarpatskyy@urk.edu.pl}
	
	%\author[1,2]{\fnm{Third} \sur{Author}}\email{iiiauthor@gmail.com}
	%\equalcont{These authors contributed equally to this work.}
	
%		\equalcont{These authors contributed equally to this work.}

	\affil[1]{\orgdiv{Mathematical Institute}, \orgname{Cardiff University}, \orgaddress{\street{Senghennydd Rd.}, \city{Cardiff}, \postcode{CF24 4AG}, 
		%\state{State}, 
		\country{UK}}}
		
	\affil*[2]{\orgdiv{Department od Applied Mathematics}, \orgname{University of Agriculture in Krakow}, \orgaddress{\street{Aleja Mickiewicza 21}, \city{Krakow}, \postcode{31-120}, %\state{State}, 
			\country{Poland}}}
	
	\affil[3]{%\orgdiv{Department}, 
		\orgname{Institute of Mathematics of NAS of Ukraine}, \orgaddress{\street{3 Tereschenkivska st.}, \city{Kyiv-4}, \postcode{01024}, %\state{State}, 
			\country{Ukraine}}}
	

	
	%%==================================%%
	%% Sample for unstructured abstract %%
	%%==================================%%
	
	\abstract{
		This paper studies a Grothendieck‑type finite‑dimensionality problem for closed linear subspaces embedded in Banach spaces. Let $S_{p}^{(q)} \subset L_{p}(M,d\mu)$ be a closed linear subspace of the Banach space $L_{p}(M,d\mu)$ defined with respect to a probability measure $d\mu$ on $M$. We prove that if $S_{p}^{(q)}$ is continuously embedded into $L_{q}(M,d\mu)$ for $q>p$, then its dimension $\dim S_{p}^{(q)} = N \in \mathbb{N}$ satisfies the estimate
		$\frac{1}{N}\left( \frac{\sqrt{\pi }%
			\Gamma (\frac{N+\tilde{q}}{2})}{\Gamma (\frac{\tilde{q}+1}{2})\Gamma (\frac{N%
			}{2})}\right) ^{2/\tilde{q}}\leq K_{p,q(m)}^{2},$ where $1/\tilde q + 1/q = 1$, $q = 2 + (p-2)2^{m} > p$ with $p \ne 2$ and $m \in \mathbb{N}$, and $K_{p,q(m)}>0$ is a bounded constant. 
		We also prove that certain closed linear subspaces of $L_{p}(M,d\mu)$ 
		consisting of continuous functions on $M$ must be finite dimensional.}
	
	%%================================%%
	%% Sample for structured abstract %%
	%%================================%%
	
	% \abstract{\textbf{Purpose:} The abstract serves both as a general introduction to the topic and as a brief, non-technical summary of the main results and their implications. The abstract must not include subheadings (unless expressly permitted in the journal's Instructions to Authors), equations or citations. As a guide the abstract should not exceed 200 words. Most journals do not set a hard limit however authors are advised to check the author instructions for the journal they are submitting to.
		% 
		% \textbf{Methods:} The abstract serves both as a general introduction to the topic and as a brief, non-technical summary of the main results and their implications. The abstract must not include subheadings (unless expressly permitted in the journal's Instructions to Authors), equations or citations. As a guide the abstract should not exceed 200 words. Most journals do not set a hard limit however authors are advised to check the author instructions for the journal they are submitting to.
		% 
		% \textbf{Results:} The abstract serves both as a general introduction to the topic and as a brief, non-technical summary of the main results and their implications. The abstract must not include subheadings (unless expressly permitted in the journal's Instructions to Authors), equations or citations. As a guide the abstract should not exceed 200 words. Most journals do not set a hard limit however authors are advised to check the author instructions for the journal they are submitting to.
		% 
		% \textbf{Conclusion:} The abstract serves both as a general introduction to the topic and as a brief, non-technical summary of the main results and their implications. The abstract must not include subheadings (unless expressly permitted in the journal's Instructions to Authors), equations or citations. As a guide the abstract should not exceed 200 words. Most journals do not set a hard limit however authors are advised to check the author instructions for the journal they are submitting to.}
	
	\keywords{closed Banach subspace, isometry, embedding, finite-dimensionality, probabilistic measure}
	
	%%\pacs[JEL Classification]{D8, H51}
	
	\pacs[MSC Classification]{46E30, 46B20, 46B25}
	
	\maketitle
	
\section{Introduction}

The problem of estimating the dimension of closed linear subspaces of the 
Banach space $L_{p}(M,d\mu )$ for $p>1$ with $p\neq 2$ is
classical in Banach-space theory. Such estimates play an important role in
operator theory, approximation theory, and related areas 
\cite{PrBl,BuBe,KaPe,KrPeSe,LySh,MiGr,Nowa,Plic}. Applications also arise in
dynamical systems and other branches of analysis 
\cite{FoSeTe,Kato,Lady-1,Lady-2,Nino,Tao,Tema}. A well-known result in this
direction is the classical theorem of Grothendieck, which provides an
estimate for the dimension of a closed linear subspace 
$S_{p}^{(\infty	)}\subset $ $L_{p}(M,d\mu )\hookrightarrow L_{\infty }(M,d\mu ),$ and its
generalization on the case of a linear closed subspace $S_{p}^{(c)}\subset
C(M,\mathbb{R})$ of continuous functions, embedded into $%
L_{2}(M,d\mu )$. In this paper we consider the related problem of estimation
of the dimension of closed linear subspaces {\ }$S_{p}^{(q)}\subset
L_{p}(M,d\mu )$ of the Banach space $L_{p}(M,d\mu )$ with
respect to a probability measure $d\mu $ on $M$, embedded into $%
L_{q}(M,d\mu )$, where $q=2+(p-2)2^{m}\ >p>1(\neq 2)$. A related version of
this problem was previously studied in \cite{BaPr-2}, however, several gaps
remain in the arguments presented there. In the present work we provide a
different approach that resolves these issues and leads to a rigorous
derivation of the corresponding dimension estimates.

The main result is stated in the following theorem.

\begin{theorem}
	\label{Tm_1} Let a closed linear subspace $S_{p}^{(q)}\subset L_{p}(M,d\mu
	),p>1(\neq 2),$ be embedded into a Banach space $L_{q}(M,d\mu )\ 
	$ for $q=2+(p-2)2^{m}\ >p>$ $1(\neq 2),$ where $\ d\mu $ is a probability
	measure on $M.$ Then $\ $ the dimension $\dim S_{p}^{(q)}=N\in \mathbb{N}$
	of the closed subspace $S_{p}^{(q)}\subset L_{p}(M,d\mu )\ $ proves to
	satisfy the inequality $\frac{\sqrt{\pi }\Gamma (\frac{N+\tilde{q}}{2})}{N^{%
			\tilde{q}/2}\ \Gamma (\frac{\tilde{q}+1}{2})\Gamma (\frac{N}{2})}\leq
	K_{p,q(m)}(S)^{\tilde{q}},1/\tilde{q}+1/q=1,$ for the bounded constant $%
	K_{p,q(m)}(S)>0,$ depending on the inclusion mapping $\
	J_{p}^{(q)}:S_{p}^{(q)}\subset $ $L_{p}(M,d\mu )$ $\hookrightarrow
	L_{q}(M,d\mu ).$
\end{theorem}

\begin{remark}
	Taking into account the estimation of the dimension $\dim S_{p}^{(q)}=N\in 
	\mathbb{N}$ of a linear closed subspace $S_{p}^{(q)}\subset $ $L_{p}(M,d\mu
	)\hookrightarrow L_{q}(M,d\mu )\ $\ for $q=2+(p-2)2^{m}>p>1(\neq 2),m$ $\in
	N,$ obtained in Theorem \ref{Tm_1}, it is interesting to analyse its
	interpretation and possible relationship to a known result from the book 
	\cite{Pisi} by G. Pisier, equivalently formulated below. \ 
\end{remark}

\begin{theorem}
	\label{Tm_1.1}In the space $L_{p}(0,1;d\lambda ),p\geq 1,$ there exists a
	linear infinite-diemnsional closed subspace $\mathcal{R}_{2}:=\overline{%
		\mathrm{span}_{\mathbb{R}}}\{r_{n}(t)=\mathrm{sign\sin }(2^{n}\pi t):t\in
	\lbrack 0,1],n\in \mathbb{N}\}\subset L_{2}(0,1;d\lambda ),$ consisting of
	the Rademacher orthonormal functions, which is a closed linear subspace of
	every $L_{q}(0,1;d\lambda ),1\leq q<\infty ,$ and for which the
	corresponding norms are proportional, that is for arbitrary $\ 1\leq
	q<\infty $ there exist constants $\gamma _{q}>0,$ such that for any $f\in $ $%
	\mathcal{R}_{q}:=\ \mathcal{R}_{2}\subset L_{2}(0,1)\hookrightarrow $ $%
	L_{q}(0,1;d\lambda ),q\geq 1,$ the norms $||f||_{q}=\gamma _{q}||f||_{2}.$
\end{theorem}

To analyse this theorem, recall that owing to the Kadec-Pie\l czy\'{n}ski
result  \cite{KaPe}, the Rademacher closed subspace $\mathcal{R}%
_{p}\subset L_{p}(0,1;d\lambda ),p>1,$ is complemented in $%
L_{p}(0,1;d\lambda ).$ In particular, this means that if our closed subspace 
$S_{p}^{(q)}\subset $ $L_{p}(M,d\mu )$ $\hookrightarrow L_{q}(M,d\mu )$
proves to contain the Rademacher closed subspace $\mathcal{R}_{p}\subset
L_{p}(0,1;d\lambda ),$ the dimension criterium of Theorem \ref{Tm_1} \
will \textit{a priori} give rise to the answer $\dim S_{p}^{(q)}=\infty ,$
that demonstrate reasonings below.

\begin{proof}
	Really, let us put $M=[0,1],$ $d\mu =d\lambda \ $ and consider the
	Rademacher closed subspace $\mathcal{R}_{p}\subset $ $L_{p}(0,1;d\lambda ),$
	mentioned above. \ Take into account two inequalities : the first one,
	stated by \ Khintchine: 
	\begin{equation}
		\ A_{\tilde{q}}(\mathcal{R})\leq ||\langle \xi |r\rangle ||_{\tilde{q}}\leq
		B_{\tilde{q}}(\mathcal{R})\   \label{P0}
	\end{equation}%
	for $\xi \in \mathbb{S}^{N-1},$ $r=(r_{1},r_{2},...,r_{N}),N\in \mathbb{N},$
	where $A_{\tilde{q}}(\mathcal{R})<B_{\tilde{q}}(\mathcal{R})$ are some
	constants with $\ B_{\tilde{q}}(\mathcal{R})=\sqrt{2}$ $\left( \frac{\Gamma (%
		\frac{\tilde{q}+1}{2})}{\sqrt{\pi }}\right) ^{1/\tilde{q}}$\ $,$ \ as it was
	calculated by Uffe Haagerup in 1981  \cite{Haag}, and the second one\ (%
	\ref{S2d}) at $\varphi =r,\mathcal{R}_{p}$ $\mathcal{\subseteq }$ $%
	S_{p}^{(q)},$ as it is calculated in the next Section below:%
	\begin{equation}
		\ ||\langle \xi |r\rangle ||_{\tilde{q}}^{-1}\leq K_{2,q(m)}(S).  \label{P1}
	\end{equation}%
	Then, as follows from Theorem  \ref{Tm_1}, one easily derives the
	inequalities 
	\begin{equation}
		\sqrt{\frac{2}{N}}\ \left( \frac{\Gamma (\frac{N+\tilde{q}}{2})}{\Gamma (%
			\frac{N}{2})}\right) ^{1/\tilde{q}}\leq 1\leq B_{\tilde{q}}(\mathcal{R}%
		)K_{2,q(m)}(S),  \label{P2}
	\end{equation}%
	which hold for all $q>1,\tilde{q}=q/(q-1),$ and arbitrary natural integers $%
	N\in \mathbb{N}.$ Thus this infinite dimensional closed subspace $%
	S_{q}^{(p)}\subset L_{p}(0,1;\mathbb{R}),$ embedded into $L_{q}(0,1;\mathbb{R%
	})$ for $q>p>1,$ is really infinite dimensional, confirming the statement of
	Theorem  \ref{Tm_2}, proving the theorem.
	
	Otherwise, if the chosen closed subspace $S_{q}^{(p)}\subset L_{p}(0,1;%
	\mathbb{R})$ is such that $\dim \left( S_{q}^{(p)}\cap \mathcal{R}%
	_{p}\right) <\infty ,$ the related estimation of its dimension $\dim
	S_{q}^{(p)}=N\in \mathbb{N}$ is given by the numerical inequality of Theorem
	 \ref{Tm_1}.
\end{proof}

In the special case, when $q=\infty ,$ as well as when closed linear
subspaces of the Banach space $L_{p}(0,1;d\mu ),p>1,$ consist of continuous
functions, there are stated the following Grothendieck type propositions.

\begin{proposition}
	Let a linear closed topological subspace $S_{p}^{(\infty )}\subset
	L_{p}(M,d\mu ),p>1(\neq 2),$ be embedded into a Banach space $%
	L_{\infty }(M;d\mu ),$ where $d\mu $ is a probability measure on $M.$ Then
	the dimension of the closed subspace $S_{p}^{(\infty )}\subset L_{p}(M,d\mu
	) $ proves to satisfy the inequality $\dim S_{p}^{(\infty )}=N\leq
	K_{p,\infty }^{2}$ for some bounded constant $K_{p,\infty }>0.$
\end{proposition}

\begin{proposition}
	Let $S_{p}^{(c)}\subset C([0,1];\mathbb{R})$ be a closed subspace of the
	Banach space $(C([0,1];\mathbb{R}),||\cdot ||_{\infty })$ of continuous
	functions on the interval $[0,1]\subset \mathbb{R}_{+},$ which allows the
	embedding into a Banach space $L_{p}(0,1;d\mu ),p>1,$ with respect
	to a probability measure $d\mu $ on $[0,1].$ Then the subspace $%
	S_{p}^{(c)}\subset C([0,1];\mathbb{R})$ is finite-dimensional.
\end{proposition}

\section{Embedding of closed subspaces into Banach spaces}

Below we consider a closed linear subspace $S_{p}^{(q)}\subset L_{p}(M;d\mu
),$ allowing the embedding into the Banach space $L_{q}(M;d\mu ), $
where $q>p>1.$ Then the following theorem holds.

\begin{theorem}
	\label{Tm_2} Let $S_{p}^{(q)}\subset L_{p}(M,d\mu )$ be a closed linear
	subspace with $p>1$ and $p\neq 2$. Assume that $S_{p}^{(q)}$ is
	embedded into $L_{q}(M,d\mu )$, where $q=2+$ $(p-2)2^{m}\ >p>1$ $(\neq 2)$, $%
	m\in \mathbb{N}$, and $d\mu $ is a probability measure on $M$. Then the
	dimension $\dim S_{p}^{(q)}=N$ proves to satisfy the determining inequality $%
	\frac{\sqrt{\pi }\Gamma (\frac{N+\tilde{q}}{2})}{N^{\tilde{q}/2}\ \Gamma (%
		\frac{\tilde{q}+1}{2})\Gamma (\frac{N}{2})}\leq K_{p,q(m)}(S)^{\tilde{q}},$
	where $1/\tilde{q}+1/q=1$ and $K_{p,q(m)}(S)>0$ is a bounded constant,
	depending on the inclusion mapping . 
	%	Let a linear closed topological subspace $S_{p}^{(q)}\subset
	%	L_{p}(M;d\mu ),p>1(\neq 2),$  be embedded into a Banach space $%
	%	L_{q}(M;d\mu )\ $ for $q=2+$ $(p-2)2^{m}\ >p>1$ $(\neq 2),$ where $\ d\mu \ $%
	%	\ is a probability measure $\ $on $M.$ Then $\ $ the dimension $\dim
	%	S_{p}^{(q)}=N\in \mathbb{N}$ of a closed subspace $S_{p}^{(q)}\subset
	%	L_{p}(M,d\mu )\ $ proves to satisfy the determining inequality $\ \frac{%
		%		\sqrt{\pi }\Gamma (\frac{N+\tilde{q}}{2})}{N^{\tilde{q}/2}\ \Gamma (\frac{%
			%			\tilde{q}+1}{2})\Gamma (\frac{N}{2})}\leq K_{p,q(m)}^{\tilde{q}},$ $1/\tilde{%
		%		q}+1/q=1,$  for some bounded constant $K_{p,q(m)}>0.$
\end{theorem}

Let us consider a closed linear subspace $S_{p}^{(q)}\subset L_{p}(M,d\mu )$
of the $\mathbb{ }$ Banach space $L_{p}(M,d\mu ),$ $p>1(\neq 2),$ with
respect to a probability measure on $M,$ satisfying, in addition, the
identical inclusion condition $S_{p}^{(q)}\subset \left( L_{p}(M,d\mu
);||\cdot ||_{p}\right) \hookrightarrow $ $\left( L_{q}(M,d\mu );||\cdot
||_{q}\right) \ $ for $q>p>1(\neq 2).$ In order to state Theorem \ref{Tm_2}
 we need some two lemmas.

\begin{lemma}
	\label{Lm_1} For any $q>p>1,$ there exists a bounded positive constant $%
	K_{p,q}(S)>1,$ such that 
	\begin{equation}
		||f||_{q}\leq K_{p,q}(S)\text{ }||f||_{p}  \label{S0}
	\end{equation}%
	for any $f\in S_{p}^{(q)}\subset L_{p}(M,d\mu )\hookrightarrow $ $%
	L_{q}(M,d\mu ).$
\end{lemma}

\begin{proof}
	As the linear subspace $S_{p}^{(q)}\subset L_{p}(M,d\mu ),$ embedded $%
	L_{q}(M,d\mu )$, $q>p>1$, is closed in $\ L_{p}(M,d\mu ),$ one can define
	the identity embedding mapping 
	\begin{equation}
		J_{p}^{(q)}:S_{p}^{(q)}\subset L_{p}(M,d\mu )\rightarrow L_{q}(M,d\mu ).
		\label{1a}
	\end{equation}%
	If a sequence $\{f_{n}:n\in \mathbb{N}\}$ $\subset S_{p}^{(q)}$ converges in 
	$S_{p}^{(q)}\subset L_{p}(M,d\mu )$ to an element 
	$f\rightarrow S_{p}^{(q)}\hookrightarrow $ $L_{p}(M,d\mu )$ with respect to the norm on 
	$L_{p}(M,d\mu )$ and simultaneously its image $\{J_{p}^{(q)}f_{n}:n\in 
	\mathbb{N}\}$ $\subset S_{p}^{(q)}\subset L_{q}(M,d\mu )$ converges to an
	element $g\in L_{q}(M,d\mu )\subset L_{p}(M,d\mu )$ with respect to the norm
	on$\ L_{q}(M,d\mu ),$ one can identify these limiting functions $f\sim g$ 
	almost everywhere. Really, since $(M,d\mu )\subset L_{p}(M,d\mu
	)\hookrightarrow L_{q}(M,d\mu ),q>p\ >1,$ from the estimations 
	\begin{equation}
		\begin{array}{c}
			||\ f-g||_{p}\leq ||f-f_{n}||_{p}+||g-f_{n}||_{p}\leq \\ 
			\leq ||f-f_{n}||_{p}+||\left( g-f_{n}\right) ||_{p}\leq \\ 
			\leq ||f-f_{n}||_{p}+||\left( g-J_{p}^{(q)}f_{n}\right) ||_{q}\ \overset{%
				n\rightarrow \infty }{\rightarrow }0%
		\end{array}
		\label{1b}
	\end{equation}%
	one obtains that $f\sim g$ almost everywhere and the image $J_{p}^{(q)}$ $%
	(S_{p}^{(q)})\subset L_{q}(M,d\mu )$ is closed in $L_{q}(M,d\mu ),q>p>1.$
	The latter, owing to the Banach closed graph theorem \cite%
	{Bana,BaLyMy,Lax,ReSi,Rudi}, gives rise to the existence of such a positive
	constant $K_{p,q}<\infty $ that 
	\begin{equation}
		||\ f||_{q}\leq K_{p,q}(S)\text{ }||f||_{p}  \label{2}
	\end{equation}%
	for arbitrary $f\in S_{p}^{(q)}\subset L_{p}(M,d\mu )\hookrightarrow
	L_{q}(M,d\mu ),q>p>1$. Remark also, that the following estimations 
	\begin{equation}
		||f||_{2}\leq ||\ f||_{q}=||J_{p}^{(q)}f||_{q}\leq K_{p,q}(S)\text{ }%
		||f||_{p}<\infty  \label{3}
	\end{equation}%
	hold for any $f\in S_{p}^{(q)}\subset L_{p}(M,d\mu )\ \hookrightarrow
	L_{q}(M,d\mu ),q>p>2,$ easily following from the Young inequality. %$\Box $
\end{proof}

Taking into account Lemma \ref{Lm_1}, we can formulate the next lemma, which
is in some sense the converse to the inequality \ (\ref{S0}).

\begin{lemma}
	\label{Lm_2}There exists a constant $K_{p,q(m)}(S)>0,$ such that the
	following inequality 
	\begin{equation}
		||\ f||_{q}\leq K_{p,q(m)}(S) ||f||_{2}  \label{S1}
	\end{equation}%
	holds for $f\in S_{p}^{(q)}\hookrightarrow L_{q}(M;d\mu ),$ $%
	q=(p-2)2^{m}+2>p>1(\neq 2),\ $ and arbitrary natural $m\in \mathbb{N}.$
\end{lemma}

\begin{proof}
	If $1<p\leq 2,$ from the Young inequality 
	\begin{equation}
		||f||_{p}\leq \ ||f||_{2}  \label{S1a}
	\end{equation}%
	for any $f\in S_{p}^{(q)}\subset L_{p}(M,d\mu )\hookrightarrow L_{q}(M,d\mu
	) $ one obtains \ inequality (\ref{S1}) $\ $for the bounded $\ $ $%
	K_{p,q}(S)>0. $\ If $p>2,$ then  one can make use of the following
	inequality: \ 
	\begin{equation}
		||f||_{p}\leq ||f||_{2^{m}(p-2)+2}^{\frac{\ 2^{m}(p-2)+2}{2^{m}p}}\text{ }%
		||f||_{2}^{\frac{2(2^{m}-1)\ }{2^{m}p}},  \label{S2}
	\end{equation}%
	which holds for any $f\in S_{p}^{(q)}\subset L_{p}(M,d\mu )$ \ and arbitrary
	natural $m\in \mathbb{N}.$ \ Now having put, by definition, $q=2+(p-2)2^{m}\
	>p>1(\neq 2),m\in \mathbb{N},$ the inequality (\ref{S2}) \ jointly with that
	of \ (\ref{S0}) gives rise to the searched estimation \ (\ref{S1}), \ where
	the constant $K_{p,q(m)}(S)=$ $K_{p,q}(S)^{\frac{(q-2)p}{2(q-p)}}$ $>0\ $\
	is bounded, thus proving the lemma.  %$\Box $
\end{proof}

\bigskip

\begin{proof}
	(\textit{Proof of Theorem \ref{Tm_2}). }Based on the lemmas above, one can
	proceed to proving Theorem \ref{Tm_2}. First we can observe that
	inequality (\ref{S1}) can be estimated, owing to the classical Young
	inequality, from the below as 
	\begin{equation}
		|l_{\varphi }(f)|\text{ }\leq K_{p,q(m)} ||f||_{2}  \label{S2a}
	\end{equation}%
	by means of a bounded linear functional $l_{\varphi }:(S_{p}^{(q)};||\cdot
	|_{p}|)$ $\rightarrow \mathbb{R}$ on the Banach subspace $%
	(S_{p}^{(q)};||\cdot ||_{q}),$ where $l_{\varphi }(f)=(\varphi
	|f):=\int_{M}\varphi fd\mu $ for some $\varphi \in (S_{p}^{(q)};||\cdot
	||_{q})^{\prime }\simeq $ $(S_{p}^{(q)};||\cdot ||_{\tilde{q}}),$ $1/\tilde{q%
	}+1/q=1,$ under the constraint $\ ||\varphi ||_{\tilde{q}}=1.$ Taking
	inequality \ (\ref{S2a}) and the evident embedding condition $%
	(S_{p}^{(q)};||\cdot ||_{2})\subset (S_{p}^{(q)};||\cdot ||_{\tilde{q}}),$
	one can calculate that 
	\begin{equation}
		\sup_{||f||_{2}\neq 0}\frac{|l_{\varphi }(f)|}{||f||_{2}}=||\varphi
		||_{2}\leq K_{p,q(m)}(S).  \label{S2b}
	\end{equation}%
	If \ now to choose an orthonormal basis $\ \Phi =\{\varphi _{1},\varphi
	_{2},...,\varphi _{N}\}\subset (S_{p}^{(q)};||\cdot ||_{2})\ $for some $N\in 
	\mathbb{N},$ $(\varphi _{j}|\varphi _{k})=$ $\int_{M}$ $\varphi _{j}\varphi
	_{k}d\mu =$ $\delta _{jk},||\varphi _{j}||_{2}=1,j,k=\overline{1,N},$ one
	can observe that a function $\varphi _{a}:=\langle a|\varphi \rangle _{N}=$ $%
	\sum_{j=1}^{N}a_{j}\varphi _{j}$ $\in $ $(S_{p}^{(q)};||\cdot ||_{2})$ has
	the norm 
	\begin{equation}
		||\varphi _{a}||_{2}=\left( \sum_{j=1}^{N}|a|^{2}\right) ^{1/2}=|a|_{N},
		\label{S2bb}
	\end{equation}%
	where the vector $\varphi :=(\varphi _{1},\varphi _{2},...,\varphi
	_{N})^{\intercal }\in \left( S_{p}^{(q)}\right) ^{N}$ and took $a\in \mathbb{%
		\ E}^{N},$ as an arbitrary vector. Having substituted the value of the norm
	\ ( \ref{S2bb}) into \ (\ref{S2b}), one obtains the inequality 
	\begin{equation}
		|a|_{N}\text{ }\leq K_{p,q(m)}(S),  \label{S3}
	\end{equation}%
	which should be combined with the imposed above condition $||\varphi _{a}||_{%
		\tilde{q}}=1.$ Taking into account that 
	\begin{equation}
		||\varphi _{a}||_{\tilde{q}}=||\langle a|\varphi \rangle _{N}||_{\tilde{q}%
		}=|a|_{N}||\langle \xi |\varphi \rangle _{N}||_{\tilde{q}}\text{ }=1,
		\label{S2c}
	\end{equation}%
	where $a\in \mathbb{E}^{N}\ $\ and $\xi :=a/|a|_{N}$ $\in \mathbb{S}%
	^{N-1},|\xi |_{N}=1,\ $ yielding jointly with \ (\ref{S3}) the inequality 
	\begin{equation}
		||\langle \xi |\varphi \rangle _{N}||_{\tilde{q}}^{-1}\leq K_{p,q(m)}(S),
		\label{S2d}
	\end{equation}%
	we can get rid of the spherical variables $\xi \in $ $\mathbb{S}^{N-1},$ if
	to apply \ to the norm equality (\ref{S2c} ) the averaging method \cite{John}
	over the unit sphere $\mathbb{S}^{N-1}. $Namely, by integrating it with
	respect to the spherical measure $d\omega _{N-1}(\xi ),\xi \in \mathbb{S}%
	^{N-1}:$ 
	\begin{equation}
		\begin{array}{c}
			|a|_{N}^{\tilde{q}}\int_{\mathbb{S}^{N-1}}d\omega _{N-1}(\xi
			)\int_{M}|\langle \xi |\varphi \rangle _{N}|^{\tilde{q}}d\mu = \\ 
			\\ 
			=|a|_{N}^{\tilde{q}}\text{ }||\text{ }|\varphi |_{N}||_{\tilde{q}}^{\tilde{q}%
			}\frac{2\sqrt{\pi ^{N-1}}\Gamma (\frac{\tilde{q}+1}{2})\Gamma (\frac{N}{2})}{%
				\Gamma (\frac{N+\tilde{q}}{2})}=\ \omega _{N-1},%
		\end{array}
		\label{S3a}
	\end{equation}%
	we can equivalently obtain from \ (\ref{S3a}) that  
	\begin{equation}
		\text{ }|a|_{N}=\frac{1}{||\text{ }|\varphi |_{N}||_{\tilde{q}}}\left( \frac{%
			\sqrt{\pi }\Gamma (\frac{N+\tilde{q}}{2})}{\ \Gamma (\frac{\tilde{q}+1}{2}%
			)\Gamma (\frac{N}{2})}\right) ^{1/\tilde{q}}  \label{A0}
	\end{equation}%
	where we denoted by $\omega _{N-1}=\frac{2\sqrt{\pi ^{N}}}{\Gamma (N/2)}$
	the surface of the $(N-1)-$ dimensional sphere $\mathbb{S}^{N-1}.$ $\ $Since
	the norm $||$ $|\varphi |_{N}||_{\tilde{q}}\leq ||$ $|\varphi
	|_{N}||_{2}=\left( \int_{M}\langle \varphi |\varphi \rangle _{N}d\mu \right)
	^{1/2}=N^{1/2},$ \ the equality \ (\ref{A0}) jointly with the condition \ (%
	\ref{S3}) yields the final \ numerical estimation 
	\begin{equation}
		\ \frac{1}{N}\left( \frac{\sqrt{\pi }\Gamma (\frac{N+\tilde{q}}{2})}{\Gamma (%
			\frac{\tilde{q}+1}{2})\Gamma (\frac{N}{2})}\right) ^{2/\tilde{q}}\leq
		K_{p,q(m)}(S)^{2},  \label{A0a}
	\end{equation}
	
	whose left hand side is bounded for those integers $N\in \mathbb{N},$ \
	which ensure the embedded subspace $S_{p}^{(q)}\subset L_{p}(M;d\mu
	)\hookrightarrow $ $L_{q}(M;d\mu ) $ at given $\tilde{q}=$ $q/(q-1),$  
	$q=2+(p-2)2^{m}\ $\ $>$ $p(\neq 2),\ m\in \mathbb{N},$ \ to be finite
	dimensional, that is $\ \mathrm{dim}$ $S_{p}^{(q)}=N<\infty .$
	
	Regarding the critical case $q=\infty $, since the representation (\ref{S2a}%
	) is not more acceptable, we need to consider that the linear bounded
	functional used there should be replaced by the following natural
	expression: 
	\begin{equation}
		|l_{x}(f)|\text{ }\leq K_{p,\infty }(S)||f||_{2}  \label{B1}
	\end{equation}%
	for any $f\in S_{p}^{(\infty )}\subset L_{\infty }(M;d\mu )\ $and $\ x\in M,$
	where the value $l_{x}(f):=f(x)\in \mathbb{R}.$ Having calculated the value%
	\begin{equation}
		\sup_{||f||_{2}\neq 0}\frac{|l_{x}(f)|}{||f||_{2}}=||l_{x}||\text{ }\leq
		K_{p,\infty }(S)\   \label{B2}
	\end{equation}%
	and using the Riesz representation theorem for the functional $%
	l_{x}:(S_{p}^{(\infty )};||\cdot ||_{2})\rightarrow \mathbb{R}$ on the
	Hilbert subspace $(S_{p}^{(\infty )};||\cdot ||_{2})\subset (L_{2}(M;d\mu
	);||\cdot ||_{2}),$ there exists for any $x\in M$ such a function $g_{x}\in $
	$(S_{p}^{(\infty )};||\cdot ||_{2})$ \ that $l_{x}(f)=(g_{x}|f)\ $and $%
	||l_{x}||$ $=||g_{x}||_{2}$ for all $f\in (S_{p}^{(\infty )};||\cdot
	||_{2}). $ If now $\Phi _{p}^{(\infty )}:=$ $\{\varphi _{1},\varphi
	_{2},...,\varphi _{N},...\}$ $\subset (S_{p}^{(\infty )};||\cdot ||_{2})\ $
	is a complete orthonormal set of functions, that is $||\varphi
	_{j}||_{2}=1,(\varphi _{j}|\varphi _{k})=$ $\int_{M}$ $\varphi _{j}\varphi
	_{k}d\mu =$ $\delta _{jk},j,k\in \mathbb{N},$ the related Parceval equality 
	\begin{equation}
		||g_{x}||_{2}^{2}=\sum_{j\in \mathbb{N}}|(g_{x}|\varphi
		_{j})|^{2}=\sum_{j\in \mathbb{N}}|\varphi _{j}(x)|^{2}  \label{B3}
	\end{equation}%
	combined with the inequality (\ref{B2}) yields the next one: 
	\begin{equation}
		\sum_{j\in \mathbb{N}}|\varphi _{j}(x)|^{2}\leq K_{p,\infty }^{2}(S),
		\label{B4}
	\end{equation}%
	which holds for any $x\in M.$ Having integrated the obtained inequality 
	( \ref{B4}) over the whole space $M,$ we obtain that 
	\begin{equation}
		\mathrm{card}\text{ }\Phi _{p}^{(\infty )}=N\leq K_{p,q}^{2}(S)  \label{B5}
	\end{equation}%
	for some $N=\dim S_{p}^{(\infty )},$ thus proving the theorem. %$\Box $
\end{proof}

The last reasonings above, concerning the special case $q=\infty ,$ can be
reformulated as the following proposition.

\begin{proposition}
	Let a linear closed topological subspace $S_{p}^{(\infty )}\subset
	L_{p}(M;d\mu ),p>1(\neq 2),$ be embedded into a Banach space $%
	L_{\infty }(M;d\mu ),$ where $d\mu $ is a probability measure on $M.$
	\ Then $\ $ the dimension of the closed subspace $S_{p}^{(\infty )}\subset
	L_{p}(M,d\mu )\ $ proves to satisfy the inequality $\dim S_{p}^{(\infty
		)}=N\leq K_{p,\infty }(S)^{2}$  for some bounded constant $K_{p,\infty
	}>0.$
\end{proposition}

Moreover, as a technical consequence of the results above the following
Grothendieck type \cite{Grot} proposition holds.

\begin{proposition}
	Let $S_{p}^{(c)}\subset C([0,1];\mathbb{R})$ be a closed subspace of the $%
	\mathbb{ }$ Banach space $L_{p}(0,1;d\mu ),p>1,$ with respect to a
	probability measure $d\mu $ on $[0,1].$ Then the subspace $%
	S_{p}^{(c)}\subset $ $C([0,1];\mathbb{R})$ is finite-dimensional.
\end{proposition}

\begin{proof}
	As a closed subspace $S_{p}^{(c)}\subset L_{p}(0,1;d\mu )\hookrightarrow
	C([0,1],\mathbb{R})$ of continuous functions on the interval $[0,1]$ can be
	closely embedded into the Banach space $C([0,1],\mathbb{R}),$
	from the Banach closed mapping theorem \cite{Lax,Rudi} one derives the
	existence of such a constant $K_{q}>1$ that the corresponding embedding
	operator $J_{p}^{(c)}:S_{p}^{(c)}\subset L_{p}(0,1;d\mu )\hookrightarrow
	C([0,1],\mathbb{R})$ is bounded, that is 
	\begin{equation}
		||J_{p}^{(c)}f||_{\infty }\leq K_{p,c}(S)||f||_{p}  \label{S14}
	\end{equation}%
	for any $f\in S_{p}^{(c)}\subset (C([0,1];\mathbb{R}),||\cdot ||_{\infty }).$
	Remark now that if $1<p\leq 2,$ then the inequality $||f||_{p}\leq ||f||_{2}$
	holds for all $f\in S_{p}^{(c)}\subset C([0,1];\mathbb{R})\hookrightarrow
	L_{p}(0,1;d\mu ).$ If $p>2,$ we can observe that $|f|^{p}\leq ||f||_{\infty
	}^{p-2}|f|^{2}$ for any $f\in S_{p}^{(c)}\subset C([0,1];\mathbb{R}%
	)\hookrightarrow L_{p}(0,1;d\mu ),$ whence by integration over the interval $%
	[0,1]$ one easily obtains that $||f||_{p}\leq ||f||_{\infty }^{\frac{p-2}{p}%
	}||f||_{2}^{\frac{2}{p}}.$ Substituting the latter inequality into  
	(\ref{S14}),  we obtain that $\ ||f||_{\infty }\leq
	K_{q,c}(S)^{p/2}||f||_{2},$ what jointly with the evident inequality $%
	||f||_{2}\leq ||f||_{\infty }$ gives rise to the dual inequality 
	\begin{equation}
		||f||_{2}\leq ||f||_{\infty }\leq K_{p,c}(S)^{p/2}||f||_{2}  \label{S15}
	\end{equation}%
	for all $f\in S_{p}^{(c)}\subset L_{p}(0,1;d\mu )\hookrightarrow C([0,1];%
	\mathbb{R}).$ \ As above, define for any $t\in \lbrack 0,1]$ a bounded
	linear functional $l_{t}:$ $(S_{p}^{(\infty )};||\cdot ||_{2})$ $\rightarrow 
	\mathbb{R}$ on the Hilbert subspace $(S_{p}^{(c)};||\cdot ||_{2}),$ such $\ $
	that $l_{t}(f)=f(t)\in \mathbb{R},$ which allows owing to the Riesz theorem
	the representation $l_{t}(f)=(g_{t}|f)$ for all $f\in (S_{p}^{(\infty
		)};||\cdot ||_{2}),$ where $g_{t}\in $ $(S_{p}^{(\infty )};||\cdot ||_{2})$
	and \ $||l_{t}||$ $=||g_{t}||_{2}.$ If now $\Phi _{p}^{(c)}:=\{\varphi
	_{1},\varphi _{2},...,\varphi _{N},...\}$ $\subset (S_{p}^{(c)};||\cdot
	||_{2})\ $ is a complete orthonormal set of functions, that is $||\varphi
	_{j}||_{2}=1,(\varphi _{j}|\varphi _{k})=$ $\int_{M}$ $\varphi _{j}\varphi
	_{k}d\mu =$ $\delta _{jk},j,k\in \mathbb{N},$ the related Parceval equality 
	\begin{equation}
		||g_{t}||_{2}^{2}=\sum_{j\in \mathbb{N}}|(g_{t}|\varphi
		_{j})|^{2}=\sum_{j\in \mathbb{N}}|\varphi _{j}(t)|^{2}  \label{S15aa}
	\end{equation}%
	jointly with the inequality \ (\ref{B2}) gives rise to the inequality 
	\begin{equation}
		\sum_{j\in \mathbb{N}}|\varphi _{j}(t)|^{2}\leq K_{p,c}(S),  \label{S15ab}
	\end{equation}%
	which holds for any $t\in M.$ Integration of the obtained above inequality \
	(\ref{S15ab}) over the whole interval $[0,1]$ yields the constraint 
	\begin{equation}
		\mathrm{card}\text{ }\Phi _{p}^{(c)}=N\leq K_{p,c}(S)^{2}  \label{S15bb}
	\end{equation}%
	for some $N=\dim S_{p}^{(c)},$ thus proving the proposition. %$\Box $
\end{proof}

\section{Conclusion}

We have closed linear subspaces $S_{p}^{(q)}$ of the 
Banach space $ (L_{p}(M,d\mu );||\cdot ||_{p}),p>1(\neq 2), $
allowing the embedding into the Banach space $(L_{q}(M,d\mu
);||\cdot ||_{q}),q>p>1(\neq 2),$ regarding a probability measure $d\mu $ on 
$M.$\ We derived the nuperical estimation \ $\ \frac{1}{N}\left( \frac{\sqrt{%
		\pi }\Gamma (\frac{N+\tilde{q}}{2})}{\Gamma (\frac{\tilde{q}+1}{2})\Gamma (%
	\frac{N}{2})}\right) ^{2/\tilde{q}}\leq K_{p,q(m)}(S),$ on the dimension $%
\dim S_{p}^{(q)}=N\in \mathbb{N}$ of a closed embedded subspaces 
$S_{p}^{(q)}\subset (L_{p}(M,d\mu );||\cdot ||_{p})\hookrightarrow
(L_{q}(M,d\mu );||\cdot ||_{q})$ $\ $into $(L_{q}(M,d\mu );||\cdot ||_{q}),$
if $q=2+(p-2)2^{m}$ $>$ $p>1(\neq 2),m\in \mathbb{N}.$  In case of the
space $M=[0,1]\subset \mathbb{R}_{+},$ endowed with an arbitrary probability
measure $d\mu ,$ we stated the Grothendieck type finite-dimensionality
result for a linear closed subspace $S_{p}^{(c)}\subset $ $\left(
L_{p}(0,1;d\mu );||\cdot ||_{p}\right) \hookrightarrow $ $\left( C([0,1];%
\mathbb{R});||\cdot ||_{\infty }\right) ,$ identical inclusion into the
Banach space $\left( C([0,1],\mathbb{R});||\cdot ||_{\infty }\right) .$ \ A
general question about estimation of the dimension of a linear closed
subspace $S_{p}^{(q)}\subset L_{p}(M,d\mu )\hookrightarrow $ $L_{q}(M,d\mu )$
for arbitrary $q>p>1$ looks to be still open and needs more sophisticated
techniques, mainly based on analysis of the complementary subspaces in $%
L_{p}(M,d\mu )$ and $L_{q}(M,d\mu ).$

\section{Acknowledgements}

The authors are much appreciated to participants of the Seminar at the
Department of Applied Mathematics, University of Agriculture in Krakow for
useful remarks, comments and suggestions. Special thanks belong to our
colleagues M. Ptak and A.K. Prykarpatski for permanent interest in our
research endeavor, fruitful discussions and instrumental help during
preparation the manuscript.

\section*{Declarations}

%	Some journals require declarations to be submitted in a standardised format. Please check the Instructions for Authors of the journal to which you are submitting to see if you need to complete this section. If yes, your manuscript must contain the following sections under the heading `Declarations':


\begin{itemize}
	\item Funding: This research was financed from the subsidy of the 
	Ministry of Science and Higher Education for the Hugo Kołłątaj 
	Agricultural University in Kraków for the year 2026
	\item Conflict of interest/Competing interests: The authors declare that they have no known competing financial interests or personal relationships that could have appeared to influence the work reported in this paper.
	\item Ethics approval and consent to participate: Not applicable
	\item Consent for publication: Not applicable
	\item Data availability: Not applicable
	\item Materials availability: Not applicable
	\item Code availability: Not applicable
	\item Author contribution: The authors contributed equally to this work.
\end{itemize}

%\begin{thebibliography}{99}
%	\bibitem{BlPr}  {\normalsize Prykarpatsky  A.K., Blackmore  D.,  A solution
%		set analysis of a nonlinear operator equation using a Leray--Sschauder type
%		fixed point approach. Topology 48,  (2009), \ p. 182--185}
%	
%	\bibitem{AlOd} {\normalsize Alspach D.E., Odell E.W., Functional Analysis:
%		Proceedings of the Seminar at the University of Texas, Editors: E.W. Odell
%		and P. Rosenthal, \ Austim, TX, USA, 1986-87 }
%	
%	\bibitem{Alze} {\normalsize Alzer H., Inequalities for the Gamma-function.
%		Proceedings of the AMS, 128(1) 1999, 141-147}
%	
%	\bibitem{AsMa-1} {\normalsize Astashkin S.V. and Maligranda L., $%
%		L_{p}+L_{\infty }$ and $L_{p}\cap L_{\infty }$ are not isomorphic for all $%
%		1\leq p<\infty ,p\neq 2.$ Proc. Amer. Math. Soc. 146 (2018), 2181-2194 }
%	
%	\bibitem{AsMa-2} {\normalsize Astashkin S.V. and Maligranda L., Rademacher
%		functions in Ces`aro type spaces, Studia Mathematica, 198(3) (2010), 235-247}
%	
%	\bibitem{BaPr-2} {\normalsize Balinsky, A.A., Prykarpatski, A.K.: On the
%		finite dimensionality of closed subspaces in }${\normalsize L}_{%
%		{\normalsize p}}{\normalsize (M,d\mu )\cap L}_{{\normalsize q}}{\normalsize %
%		(M,d\mu ).}$ \ {\normalsize Axioms, 10(2022), \ p. 275 }
%	
%	\bibitem{Bana} {\normalsize Banach S., Theorie des operations lineaires,
%		PWN, Warszawa, 1932. }
%	
%	\bibitem{BaBaPlPr} {\normalsize Banakh I., Banakh T., Plichko A.,
%		Prykarpatsky\ A.K., On local convexity of nonlinear mappings between Banach
%		spaces. Cent. Eur. J. Math., 10(6), (2012), p. 2264-2271 }
%	
%	\bibitem{BaLyMy} {\normalsize Banakh T., Lyantse W.E. , Mykytyuk Ya.V., $%
%		\infty $-Convex sets and their applications to the proof of certain
%		classical theoremsof functional analysis. Matematychni Studii. 11(1),
%		(1999), p. 83-84 }
%	
%	\bibitem{BuBe} {\normalsize Butzer P.L., Berens H., Approximation,
%		Springer-Verlag Berlin Heidelberg, 1967 }
%	
%	\bibitem{FoSeTe} {\normalsize Foias C., Sell G.R., Temam R. Inertial
%		mathifolds for nonlinear evolution equations. J. Diff. Eqns., 73, (1988), p.
%		309-353 }
%	
%	\bibitem{Grot} {\normalsize Grothendieck A., Sur certains sous-espaces
%		vectoriels de $L_{p}.$ Canadian J. Math. 6 (1954), p. 158--160 }
%	
%	\bibitem{John} {\normalsize John F., \ Plane waves and spherical means.
%		Interscience Publisher, London, 1955}
%	
%	\bibitem{JoSh} {\normalsize Johnson W.B. , Schechtman G., Finite dimensional
%		subspaces of $L_{p}.$ Elsevier Preprint, January 31, 2001, p. 1-40 }
%	
%	
%	\bibitem{Haag} {\normalsize Haagerup U., The best constants in Khintchine
%		inequality, Studia Mathematica, LXX (1982), p. 232-283}
%	
%	\bibitem{KaPe} {\normalsize Kadec M.I, Pe\l czy\'{n}ski A., Bases, lacunary
%		sequences and complemented subspaces in the spaces }$L_{p}${\normalsize ,
%		Studia Math., 21 (1962), 161-176}
%	
%	\bibitem{Kato} {\normalsize Kato T. Nonlinear evolution equations in Banach
%		spaces, Proceed. Symp. On Pure Appl. Math., 45(2), (1986), p. 9-23 }
%	
%	\bibitem{KrPeSe} {\normalsize Kre\u{\i}n S.G. , Petun\={\i}n Yu.I, and Sem%
%		\"{e}nov E.M., Interpolation of linear operators, Translations of
%		Mathematical Monographs, vol. 54, American Mathematical Society, Providence,
%		R.I., 1982. }
%	
%	\bibitem{Lady-1} {\normalsize Ladyzhenskaya O.A., Finite-dimensionality of
%		bounded invariant sets for Navier--Stokes systems and other dissipative
%		systems, Zap. Nauchn. Sem. LOMI, 1987, 163, 105--129 }
%	
%	\bibitem{Lady-2} {\normalsize Ladyzhenskaya O.A., Estimates of the fractal
%		dimension and of the number of deterministic modes for invariant sets of
%		dynamical sytsems, Zap. Nauchn. Sem. LOMI, 1987, 163, p. 105--129;
%		https://www.researchgate.net/publication/2428834; Finite dimensional
%		subspaces of $L_{p}.$ }
%	
%	\bibitem{Lax} {\normalsize Lax P.D., Functional analysis, Wiley-Interscience
%		2002.}
%	
%	\bibitem{Lewi} {\normalsize Lewis D.R., Ellipsoids defined by Banach ideal
%		norms, Mathematika, 26(1), (1979), p. 18-29 }
%	
%	\bibitem{LiLo} {\normalsize Lieb E.H., Loss M., Analysis, Graduate Studies
%		in Mathematics, v. 14, American Mathematical Society, Providence, R.I., 1997
%	}
%	
%	\bibitem{LiWu} {\normalsize Liu Y., Wu J. The Grothendieck theorem in
%		Bergman spaces. arXiv:2511.01521v1 [math.CV] 3 Nov 2025 }
%	
%	\bibitem{LySh} {\normalsize Lyubich Yu.I., Shatalova O.A., Isometric
%		embeddings of finite-dimensional }$l_{p}${\normalsize -spaces over the
%		quaternions, Algebra and Analysis, 16(1) (2005), 9--24}
%	
%	\bibitem{MiGr} {\normalsize Michalak A., Grala-Michalak J.,On constructions
%		of isometric embeddings of nonseparable }$L_{p}${\normalsize \ spaces, }$%
%	0<p\leq 2.${\normalsize \ Commentationes Mathematicae, 48(2) (2008), 129-145}
%	
%	\bibitem{Nino} {\normalsize Ninomiya H. Some remarks on inertial manifolds.
%		J Math.Kyoto Univ., 32(4) 1992, p. 667-688 }
%	
%	\bibitem{Pisi} {\normalsize Pisier G., \ The Volume of Convex Bodies and
%		Banach Space Geometry, Cambridge University Press, 2010}
%	
%	\bibitem{Plic} {\normalsize Plichko A., Rate of Decay of the Bernstein
%		Numbers, Journal of Mathematical Physics, Analysis, Geometry} {\normalsize %
%		9(1) 2012, 1-14}
%	
%	\bibitem{PrBl} {\normalsize Prykarpatsky A.K., Blackmore D. A solution set
%		analysis of a nonlinear operator equation using a Leray Schauder type fixed
%		point approach. Topology 48 (2009), p. 182-185 }
%	
%	\bibitem{ReSi} {\normalsize Reed M. and Simon B., Methods of modern
%		mathematical physics, vol. 1. Functional analysis, Academic Press, New York
%		1972.}
%	
%	\bibitem{RoSe} {\normalsize Rodin V.A. and Semyonov E.M., Rademacher
%		series in symmetric spaces, Anal. Math. 1 (1975), 207--222}
%	
%	\bibitem{Rudi} {\normalsize Rudin W., Functional Analysis, McGraw-Hill, 1991.%
%	}
%	
%	\bibitem{Tao} {\normalsize Tao T. An introduction to measure theory,
%		Graduate Studies in Mathematics, v. 126, American Mathematical Society,
%		Providence, R.I., 2011 }
%	
%	\bibitem{Tema} {\normalsize Temam R. Infinite-dimernsional dynamical systems
%		in fluid mechanics, Proceed. Symp. On PureAppl. Math., 45(2), (1986), p.
%		431-445 }
%	
%	\bibitem{Nowa} {\normalsize Nowak P.W., On coarse embeddibility into }$l_{p}$%
%	{\normalsize  spaces and conjecture of Dranishnikov, arXiv:math/0410566v1
%		[math.MG] 27 Oct 2004}
%\end{thebibliography}
	
	% Loading bibliography database
	\bibliography{manuscript-refs}
	
\end{document}
